\section{Stream Management Service}

\subsection{Overview}

In the SNS DAS system, there are many sources of information. Neutron data is
generated from one or more detectors for each beamline, as well as from at
least one beam monitor. Information about the evironmental conditions and
positioning of the sample under study is generally provided by ``Slow
Controls'', though certain pieces -- such as strain force and chopper phase
information -- come in as ``Fast Metadata''.  All of this information must be
connected to user-provided information such as the sample material being
investigated and which proposal number for which the experimental data is being
generated.

There are currently several consumers of the totality of this information, and
it makes little sense to have each one of these applications need to speak
to multiple data sources, each often using different protocols to obtain the
data. It becomes much more simple to have a single clearing house to which
interested parties can go to obtain data relating to the current experiment --
an ``oracle'', if you will.

The Stream Management Service (SMS) provides this functionality for the DAS
system. It will operate as a Unix-style daemon, and aggregate experiment data
into a coherent data stream. It will further manage certain pieces of
experiment metadata such as the ``run number''. While the SMS is not
user-facing, it will provide an interface for those applications which are
to associate user-collected meta-data with the appropriate experimental data.

\decision{Beam Monitor Input}{We need to decide how to bring beam monitor data
into the new stream. Currently monitor data is in a histogram format, but it
would be preferable to bring it in as another detector input.}

\todo{Glossary}{We need a glossary or other section in which we
introduce EPICs process variables and other terms used throughout this
document.}
\todo{Resetting environment to defaults}{There are a number of configuration
items that may be set by the user via EPICS PVs. It seems likely that we will
want a method to reset those to default values to make things easy during
changeovers.}

\subsection{Detector Input}

\assumption{Preprocessor Rework}{There is a project to change how the
detectors and Data Stream Packetizers (DSPs) communicate with the rest of
the DAS system. This project is not currently considered in this document.
We assume that the data will be made available in the new protocol format,
either as the result of a rework of the DSP/preprocessor interface, or via
a service that translates the existing network protocol to the new format.}

Neutron time-of-flight data and ``Fast Metadata'' are currently delivered
through the ``preprocessor'', a process that communicates with the detector
hardware and translates that data to the streaming protocol.

\subsubsection{Detector (Preprocessor) Communications}

Upon startup of the SMS daemon, a configuration file will be parsed to
determine how many preprocessors are in the system for this beamline. Each
preprocess will be identified by an address and TCP port, and other information
will be required as deemed appropriate. Each preprocessor will be identified by
a set of EPICs process variables. Many of these will be read-only and provide
access to the current configuration loaded into SMS as well as the connectivity
and status information. Each preprocessor will have the ability to be
configured into ``bypass'' or ``discard'' modes via EPICs.  These modes will
control which event sources are required for a given run, and how to handle
communication disruptions.

The SMS will attempt to maintain communications at all times with each
preprocessor configured. The SMS will initiate the TCP connection to the
preprocessor at the address and port given in the config file. Once connected,
the preprocessor will stream its data in the new format to the SMS until the
connection is closed. Should the connection be broken, whether by a timeout,
remote service reset, or deliberate acction on behalf of the preprocessor, a
new connection attempt will be immeditately scheduled. Should this connection
attempt fail (timeout, refusal, or otherwise), another attempt will be
scheduled for a configurable time in the future.

When connected, all data sourced by the preprocessor will be processed and
stored onto local stable storage as described in section~\ref{sec:storage}.
This processing will occur whether or not there is currently an experimental
run being recorded. Based on a combination of source and pixel id, each
packet will be passed to the appropriate transformation engine. Packets
containing accelerating timing information will be parsed and the values fed
to the appropriate tranformation engines as needed. Unknown packets will be
logged and dropped.

\assumption{Single-type Raw Packets}{It is currently assumed that a raw event
packet from the preprocessor will include neutron event data or fast metadata,
but not both at the same time. However, it should not be a huge undertaking to
change this assumption, as we could demultiplex based on pixel ids instead
of packet source.}

\decision{Hostname or IP?}{We need to decide if we will uses host names (FQDNs)
or IP addresses in the configuration file. Host names have the advantage of
human readability and can contain additional semantic information, but carry
cost of a dependency on DNS or a local host file. If host names are employed,
we will also need to decide if the lookup is done during the parsing of the
configuration file, or if it will need to be performed each time we try to
connect to a preprocessor. The latter will allow for changes without a full
restart of SMS, but add some additional complexity to the code.}
\decision{Single Purpose Preprocessors?}{Is it expected/permissible for a
single preprocessor to send both neutron and fast metadata?}
\decision{Handling Reported Errors}{There are provisions in the packet format
to report hardware errors from the detectors. This information must be
incorporated into our design.}

\todo{Configuration Format}{A format must be defined for the configuration
file.}
\todo{Preprocessor Info}{We need to determine what information is needed for
each preprocessor. Some of this will be dictated by the transformation
done on the stream.}
\todo{Preprocessor PVs}{Add a table to this document containing the PVs
associated with the preprocessor and their form and function.}

\subsubsection{Neutron Event Transformation}

As each raw event is processed, the pixel id will be remapped to its new value
given by a lookup table. This lookup will also give us the bank associated
with the pixel. Any pixel id that does not have a mapping in the table will
have its high-bit (bit 31) set to indicate a raw pixel value, and the bank will
be set to all ones ({\raise.17ex\hbox{$\scriptstyle\sim$}}0).  Once the pixel
is remapped, the new ID and time-of-flight value will be appended to a buffer
corresponding to the appropriate detector bank.

The pixel mapping and banking will be loaded from a configuration file specific
to each beamline. This mapping shall be a one-to-one, reversible mapping such
that no two input values map to the same output value. The revisibilty property
shall be enforced during load of the configuration, and violation will be
considered a fatal error preventing startup of the service.

Once the end of raw frames for a pulse has been detected, the current buffers
for each bank will be pushed to the aggregation component for storage and
publication to interested subscribers. These buffers will be in a packet format
for processed neutron event data, similar but separate from the raw event
format.

\assumption{Pixel ID High Bits}{The current system stores various pieces of
information in the high bits of each pixel. We will need to account for that
data, and extract it as part of processing. The reversible mapping algorithm
requires upstream to leave clear a known bit in order to indicate a mapping
failure. We currently have chosen bit 31, but this could be changed.}

\decision{Frame Correction}{Neutrons may/often take more than the nominal
16.67ms inter-pulse gap (IPG) to transit from the moderator to the detectors.
The length of time is determined by the settings of the chopper(s), the
geometry of the beamline, the position of the detector, and other factors.
Neutrons that take longer than the IPG will show up in the event data
one or more frames (pulses) after their originating pulse, and require
corrections to their time-of-flight. Currently this is done during analysis,
but there may be some desire to move that upstream. We need to decide if
that is the case, what would be required to do it, and how to ensure that
the analysis tools can easily tell corrected from un-corrected data.}
\decision{Handling Missing Packets}{The raw packets contain a sequence number
so that we may be able to detect dropped packets. What do we do in that case?}

\todo{Define Remapping/Banking Config Format}{This could be a straight
key-value table, though I understand there may be an existing language for
this.}
\todo{Determine How To Identify End-of-Pulse}{We need to know when we have
received all data for a given pulse so that we know when to push things
through the pipeline. It appears that the raw frames have a sequence number
and an end-of-pulse flag we can use for this purpose. We may also be able
to use the presence of raw event packets with a different pulse ID field, or
the reception of an RTDL packet to indicate EOP.}
\todo{Define Processed Neutron Event Packet}{This may already be in Lloyd and
Jim's document describing the streaming procotol. If not, it will need to
be defined -- it will be very similar to the raw event packet, but will
have bank information and other data as needed from the RTDL packets or
other sources.}

\subsubsection{Fast Metadata Transformation}

Raw event packets from fast-metadata sources will be passed here for
transformation into protocol metadata packets. The pixel ids will be matched to
variable ids loaded from the configuration file. As each pixel is translated,
the variable change packet will be forwarded to the aggregator without waiting
for the end-of-pulse.

\assumption{Frame Correction and Metadata}{Because the fast metadata
information is not based on neutron flight time, the pulse timestamp will
always be correct and no frame correction is needed.}

\decision{How to handle unknown fast-metadata pixel ids?}{The metadata format
for the new protocol expects there to be a ``Variable Descriptor'' (actual
term likely differs) in the stream prior to the variable being updated. If
we get an unexpected pixel id for fast metadata, what do we do? Make up a
descriptor and carry on? It would seem we want to record the data given, even
if it is from a glitch.}

\todo{Define Variable Config File Format}{We need to configure the mapping
from source/pixel id to metadata variable. We will need the XML descriptor
to insert into the stream for this variable, and may wish to provide for
simple transformations on the value to convert from an integer to the
desired units. It would be nice to have a format that avoids my parsing XML,
but this is not a hard requirement and XML parsing may be required for other
functionality in any event.}

\subsection{Slow Controls}

The SMS must maintain a list of descriptions and current values of the PVs
describing the sample environment. This is needed in order to provide
new aggregation clients the current operating paramters as well as to
record the environment for each experiment run. 

\decision{How to integrate slow controls?}{There have been discussions of SMS
being a direct EPICs client to obtain PVs and incorporate those into the
stream. This approach seems to be duplicative of work that will be needed in
the scan server and perhaps other parts of the system, and as such -- while not
completely off the table -- is not the currently favored approach. More
discussions are required to determine a coherent approach to slow controls for
the system.}

\todo{Scan Server Handshake}{Assuming we go forward with the scan server
providing all-things-sample-environment, we will need to determine the
steps required for the SMS and the Scan Server to establish communications.
We will also need a mechanism to indicate that the environment has changed
and to rebuild the description.}

\subsection{Run Management}

The SMS will be responsible for managing the current run identifier. It is
expected that there will be PVs for the {\bf previous}, {\bf current}, and {\bf
next} run number to be used. The PVs for current and next are mutually
exclusive, with only one being valid at a time; which one is valid will depend
on if we are currently recording an experimental run for translation. If we are
in an experimental run, then {\bf current} will contain the run number and {\bf
next} will be blank. If we are not in an experimental run, then {\bf next} will
contain the next run number to be used and {\bf current} will be blank. When
the user moves state from ``recording'' to ``end'' (our corresponding
transition would to move the ``record'' PV from true to false,) we will move
the current run number to {\bf previous} and generate a new number for {\bf
next}. When we transition from ``end'' to ``recording'' (``record'' to true) we
would move the number from {\bf next} to {\bf current}.

\addvspace{\baselineskip}
\noindent
Note: the aggregated data stream will always be sent to the local disk.
``Recording'' in this context means that we are sending this data to the
translation service.

\todo{Confusing Run-number PVs}{The current description for communicating run
numbers seems a bit confusing and could probably be simplified. It may be
better to always have all three fields populated and make sure the GUI is doing
the right thing. We could also just have two PVs, with the {\bf current} and
{\bf next} values merged; the meaning would still be determined by the
recording state, but we'd just use the one PV.}

\subsection{Stream Storage}
\label{sec:storage}

The SMS daemon will store the aggregated stream to local stable storage. The
top-level directory for this storage will be specified in the configuration
file.  The subdirectory used for the current stream will be determined by the
time (UTC) we last changed states as given by the {\tt gettimeofday()} system
call. The base subdirectory name will be ``YYYYMMDD-HHMMSS.SSSSSS''; if we are
currently recording a run, then ``-run-'' and the current run number will be
appended.

There will be a separate directory provided for storage of experimental runs
flagged as in permament error by the translation service. As this will move the
streams outside of the SMS's pervue, this area will need to be monitored and
commands provided to manually resubmit the runs to the translation service once
the offending issue has been repaired. While this mechanism is not intended to
be needed outside of the testing environment, Murphy's Law dictates that this
must be designed to be usable at 3am localtime.

On startup, the SMS daemon will generate a list of directories within the
storage area. Each entry in the list will contain the base path and the total
space consumed by the files in the directory. Those directories associated with
experimental runs will be marked as such and be placed on additional list(s)
used to manage interaction with the translation service. The experimental run
entries will have a field to indicate a successful translation.

Once the storage area is scanned upon startup, the size will be checked
against a configurable limit, and the data files will be pruned as necessary.
Files marked as runs that have not yet been successfully translated will be
exempt from pruning.

As each data packet is presented to the aggregator, it will be appended to the
current data file. The file's length will be updated, as will the total storage
used. The new data is published to subscribers by informing them of the change
to the file size. Should the length of the file exceed the configured limit on
file size, an end-of-file packet will be written to indicate that the file is
not corrupt, and a new data file will be started. This file size limit provides
the granularity at which space is reclaimed from the system, and may be tuned
for each beamline based on local characteristics.

When starting a new file, whether from hitting the file-size limit or from
starting a new experimental run, the current sample environment state must be 
put in the file, along with the experiment information, geometry, translation
tables, and any other data required to understand the neutron events contained
in the file. This replication of data into each file ensures that it is not
lost, rendering the neutron data useless, as old data files are pruned.

At the beginning of a data file, and again at configurable intervals, the SMS
storage process will add signature packets to the stream stored on disk. These
packets will be used to identify the file as belonging to the DAS system and to
provide some level of data recovery in the event that an IO error occurs during
readback of the file.

At periodic intervals based on event timestamps, the storage system will add an
entry to index for the current data file. This entry will include the current
pulse ID, the file and offset of the first packet describing data related to
the pulse ID, and a snapshot of the sample environment. This index will be used
to provide the clients of the Event Stream Server the ability to request data
from a given timestamp. It is expected that there will be limited granularity
for these entries -- on the order of 5 to 15 minutes. This value is expected to
be configured for each beamline to balance index growth against minimizing the
transfer of extra event data prior to the timestamp requested.

\decision{What do we do if we exhaust our storage?}{While this pruning is
expected to keep the storage area managable, there is the risk of exhausting
the space if the translation services are down for an extended period and/or
other data is placed on the same partition without adjusting the SMS's idea of
the maximum space it should use. Should the SMS keep track of the free space on
the disk and try to prune itself even if it is below its limit? What do we do
if we hit an ENOSPC error when writing event data? How hard do we try to
recover vs expecting the system managers to avoid reaching that point?}
\decision{What do we do on IO error?}{Disk IO is a physical process, and we are
bound to hit IO errors at some point in the life of the system. How do we plan
to handle them?}
\decision{Hash data for corruption checking?}{Do we want to include a hash of
the contents of the file to the end-of-file packet to be able to check for
corruption? This adds computational overhead, and the gain may not be worth it
for these temporary files. The translation service may or may not want this for
its archival files, but it can filter and add its own hash as it sees fit.}

\todo{Management policy for index data}{We will need to determine the pruning
policy of the timestamp index data.}
\todo{Add PVs for monitoring}{I expect we will define some PVs to export
various state information about space used, etc.}


\subsection{Translation Client}

The SMS daemon will manage the transmission of ``recorded'' experimental runs
to the translation service (TS). For each run to be translated, the SMS will
open a separate connection to the TS. The data from the file(s) associated with
that run will be relayed across the the socket until the end of the run. Once
the data has been transfered to the TS, the SMS will wait for a response packet
indicating the success or failure of the translation.  On a successful
translation, the SMS will mark the current run as translated and eligible for
pruning.

If there is a communication error or the TS indicates that an error occurred,
the run will \emph{not} be marked as translated. On receipt of a permament
error indication from TS, the run directory will be moved to the error area,
and a message logged. For transient errors, the run will placed on a list for
future translation, after a delay specified by the TS. After the delay passes,
the entire experimental run will be resent to the TS. There is will be a
configurable retry limit for transient error attempts to avoid repeatedly
attempting to translate a run with issues and impeding forward progress on
other runs.  Runs exceeding this retry limit will also be moved to the error
area. Communication errors will be handled as transient errors, with a retry
delay specified from the configuration file.

When a backlog of experimental runs are queued up for transmission to the TS,
the SMS will allow for a configurable number of connections to be opened. This
will allow a beamline to realize the benefits of having a cluster of
translation servers. One slot will be reserved for the current experimental run
on the theory that the users should continue to get new data as soon as
possible. This mechanism will need tuning to avoid having beamlines monopolize
the bandwidth/CPU/memory available to the TS, and may benefit from being
controlled via a PV to allow administrators to make on-the-fly adjustments
based on operating conditions. We will need to restrict write-access from users
-- or at least limit the range that they can set -- in order to avoid
unintentional denial of service attacks.

As the translation client will be using OS facilities to avoid memory copies,
the translation service will need to tolerate certain packets that make little
sense in a network stream context. This includes signature packets and
end-of-file packets. It is expected that the TS will discard these packets and
insert its own instance of them as needed by its local configuration aimed more
at archival than transient storage.

\decision{Translation Status Indication}{It would be good to give the user some
insight as to the progress of translation. Jim mentioned a service monitor that
exists, but there seemed to be some potential difficulties with integrating it
into the proposed systyem that I do not currently recall. As for the SMS's
side, we may want to provide basic information about if we're currently
connected, and how big our translation backlog is.}
\decision{Queue service order for translation}{We need to have a policy for
which order we send our backlog of queued runs to the TS. It seems like we may
wish to work from more recent back to older, but this could lead to those
getting farther and farther back in the queue in the face of ongoing problems.
We could over a user-settable PV with ennumerated values {\bf Recent first},
{\bf Oldest first}, and {\bf Candle}.}
\decision{Setting NeXus histogram parameters}{We need to decide how to set the
histogram parameters for use by the translation service. It seems like this
should be user controlled at the beamline, though there will be some
restrictions on what can be request, and what kind of changes, if any, would be
allowed during a run. Would these come in as regular metadata or would we give
them their own packet type? Does it go into the run info?}

\todo{Host name handling}{If load balancing for the TS is handled by DNS
round-robin, then we may need to do a DNS lookup here. If it is handled by port
mapping by an intermediate network device, we may not need to do the lookups,
but we may have some issues getting full bandwidth out of it.}


\subsection{Event Stream Server}

The SMS will provide a mechanism to allow clients to subscribe to the
aggregated event stream. The SMS will listen on a well-defined port for new
subscribers. Once connected, subscribers will send an initialization packet to
establish their requested service.

\decision{Passing run markers}{We will want to pass run recording/end events to
listeners for their perusal. While those events will generated by SMS, there
will also need to be events from the scan server to indicate points of interest
and/or when to reset counters for streaming analysis tools. We need to decide
what those look like and what their semantics will be.}
\decision{Allow metadata only clients?}{Life becomes much easier for the SMS if
we only provide access to the full aggregated stream and do not do any
filtering on behalf of subscribers. If we wish to provide a reduced stream for
allowing external monitoring or other use, we may wish to write a client that
performs that task and only forwards what it needs.}

\todo{Elaborate on Event Stream Server}{There have been many conversations
about this topic, such as allowing for a simple request to go back to a
timestamp or beginning of current run, whichever is more recent. I need to
capture those in this section, but the screen is going blurry.}


\section{Stream Histogram Tool}

This is a client of the Event Stream Server. It listens to the aggregated
stream and generates histograms based on parameters controlled via PVs. The
histogram will also be provided as the contents of a PV. Another PV will
contain the number of events contained in a region defined by yet more PVs.

\todo{Elaborate on Simple Stream Histogram}{Screen is blurry, more later.}


